%% ============================================================
%%  §10  Evaluation
%% ============================================================

\section{Evaluation}
\label{sec:evaluation}

This section reports quantitative measurements of the lean-pl-verify artifact:
theorem counts, lines of code, the end-to-end Charon pipeline,
proof-to-code ratios, build times, and a comparison with related work.
The complete project builds with \textbf{0 sorry} across all 258 theorems.

\subsection{Theorem Counts}

Table~\ref{tab:theorems} gives the complete theorem count by module.
All theorems are fully kernel-checked.
\texttt{Translation/Adequacy.lean} proves both directions of adequacy:
soundness (every interpreter output is relationally derivable) and
completeness (every relational derivation is witnessed by some fuel bound).

\begin{table}[t]
\centering
\caption{Theorem counts by module. Sorry count: \textbf{0} across all files.}
\label{tab:theorems}
\small
\begin{tabular}{@{}llrr@{}}
\toprule
\textbf{Module} & \textbf{Contents} & \textbf{Theorems} & \textbf{Sorry} \\
\midrule
\multicolumn{4}{@{}l}{\textit{Foundation}} \\
\texttt{Foundation/Monad.lean}        & Monad laws, rpanic lemmas              &  2 & 0 \\
\texttt{Foundation/Ownership.lean}    & FBPair identity laws                   &  4 & 0 \\
\midrule
\multicolumn{4}{@{}l}{\textit{Specification}} \\
\texttt{Spec/Satisfies.lean}          & ProgramSpec combinators                & 17 & 0 \\
\texttt{Spec/Examples.lean}           & RustM-level examples (E1--E8)          & 15 & 0 \\
\midrule
\multicolumn{4}{@{}l}{\textit{LLBC / Rust verification}} \\
\texttt{Translation/ElabSpec.lean}    & LLBC Rust theorems (F1--F16)           & 48 & 0 \\
\texttt{Translation/LoopInvariant.lean} & \texttt{sum\_to} loop invariant      & 14 & 0 \\
\texttt{Translation/FactInvariant.lean} & \texttt{fact} loop invariant         & 14 & 0 \\
\texttt{Translation/FibInvariant.lean}  & \texttt{fib} loop invariant          &  9 & 0 \\
\texttt{Translation/Semantics.lean}   & Relational big-step semantics          &  0$^\dagger$ & 0 \\
\texttt{Translation/Adequacy.lean}    & Full adequacy (soundness + completeness) & 18 & 0 \\
\midrule
\multicolumn{4}{@{}l}{\textit{Charon-extracted Rust (end-to-end pipeline)}} \\
\texttt{Translation/CharonDefs.lean}  & Auto-generated \texttt{LLBCFunDef}s    &  0$^\ddagger$ & 0 \\
\texttt{Translation/CharonSpec.lean}  & Proofs over Charon-generated defs      & 57 & 0 \\
\midrule
\multicolumn{4}{@{}l}{\textit{Real-crate extraction (\texttt{num-integer} v0.1.45)}} \\
\texttt{Translation/NumIntegerSpec.lean} & \texttt{Integer::is\_even<u32>} from published crate & 5 & 0 \\
\midrule
\multicolumn{4}{@{}l}{\textit{Proof automation}} \\
\texttt{Tactic/Examples.lean}         & One-liner proofs via macros            & 16 & 0 \\
\midrule
\multicolumn{4}{@{}l}{\textit{TypeScript (cross-language unification)}} \\
\texttt{TypeScript/ElabSpec.lean}     & TS theorems T1--T14 + 6 unif.\ thms   & 33 & 0 \\
\texttt{TypeScript/BugDetection.lean} & Bug detection case study (B1--B6)      &  6 & 0 \\
\midrule
\textbf{Total}                        &                                        & \textbf{258} & \textbf{0} \\
\bottomrule
\end{tabular}
\end{table}

\noindent
$^\dagger$ \texttt{Semantics.lean} defines the relational big-step semantics
and supporting lemmas; its correctness is witnessed by the adequacy theorems
in \texttt{Adequacy.lean}.

\noindent
$^\ddagger$ \texttt{CharonDefs.lean} contains 18 auto-generated
\texttt{LLBCFunDef} values (definitions, not theorems); all 57
correctness theorems reside in \texttt{CharonSpec.lean}.

\subsection{The Charon End-to-End Pipeline}
\label{sec:eval-charon}

A central engineering contribution is the pipeline that connects Rust
source code to kernel-checked Lean proofs with no manual transcription:

\begin{enumerate}[leftmargin=2em]
  \item \textbf{Rust source} (\texttt{examples/rust-crate/src/lib.rs},
        148 lines, 18 functions).
  \item \textbf{Charon} v0.1.197 (\texttt{charon cargo --dest-file ...})
        translates MIR to LLBC JSON (\texttt{verified\_fns.llbc}).
  \item \textbf{charon2lean.py} (463 lines) translates the LLBC JSON to
        \texttt{LLBCFunDef} Lean terms, handling:
        flat-list-to-CPS conversion,
        \texttt{AddChecked}/\texttt{SubChecked}/\texttt{MulChecked} desugaring,
        overflow-check \texttt{Assert} elision,
        negative integer literal wrapping,
        and dict-form binary operators (\eg \texttt{\{"Rem": "UB"\}}).
  \item \textbf{\texttt{CharonDefs.lean}} (auto-generated) is
        elaborated by Lean's kernel without error.
  \item \textbf{\texttt{CharonSpec.lean}} proves
        57 correctness theorems over the auto-generated definitions
        using the same proof tactics as \texttt{ElabSpec.lean}.
\end{enumerate}

\paragraph{Pipeline validation: adding \texttt{pow} and \texttt{gcd}.}
To validate that the pipeline is extensible beyond the initial 16 functions,
we added \texttt{pow} (integer power, loop) and \texttt{gcd}
(Euclidean algorithm using \texttt{\%}).
This required: (1)~writing the Rust functions and re-running Charon;
(2)~fixing \texttt{charon2lean.py} to handle the \texttt{\{"Rem": "UB"\}}
dict-form operator that Charon emits for signed remainder;
(3)~adding \texttt{.bitAnd} on \texttt{bool\_} to \texttt{evalBinOpPure}
(Charon's UB-check intermediate uses bool-typed bitwise-and);
(4)~writing 10 new theorems.
The proofs are all by \texttt{rfl}: \texttt{gcd(a, 0) = a} and
\texttt{pow(x, 0) = 1} hold symbolically (loop exits immediately),
and ground instances (\texttt{gcd(12, 8) = 4}, \texttt{pow(2, 10) = 1024})
reduce by direct kernel evaluation.

The pipeline demonstrates that the proof infrastructure is \emph{not tied to
hand-crafted definitions}: Charon-generated ASTs, including Charon's
UB-check temporaries, are handled uniformly by \texttt{charon\_simp}.

\subsection{Real-Crate Extraction: \texttt{num-integer} v0.1.45}
\label{sec:eval-numinteger}

To demonstrate that the framework connects to real published crates, we
ran Charon directly on the source of \texttt{num-integer} v0.1.45
(\url{https://crates.io/crates/num-integer}, $\approx$45M downloads).

\paragraph{Extraction.}
Running \texttt{charon cargo} on the crate source produced an LLBC file
with 1\,657 function entries: 359 with concrete bodies, 590 opaque
(standard library intrinsics and trait stubs).
Among the concrete bodies is the macro-generated implementation of
\texttt{Integer::is\_even} for each primitive type.
We extracted the \texttt{u32} implementation (11 LLBC statements):

\begin{lstlisting}[caption={Source: num-integer/src/lib.rs (impl\_integer\_for\_usize! macro).},
                   label={lst:iseven}]
fn is_even(&self) -> bool { *self % 2 == 0 }
\end{lstlisting}

\paragraph{Charon's output.}
Charon adds a \texttt{RemainderByZero} UB check
(\texttt{var4 = (2 == 0)}, always false) and a subsequent \texttt{Assert}
(dropped by \texttt{charon2lean.py}).
The resulting \texttt{LLBCFunDef} body is:
copy \texttt{*self} $\to$ \texttt{var3}; compute \texttt{var2 = var3 \% 2};
compute \texttt{var0 = (var2 == 0)}; return.

\paragraph{Key challenge.}
Using \texttt{num-integer} as a \emph{dependency} (the natural way)
produces only opaque bodies — Charon cannot see into compiled
dependency crates. Running Charon on the crate source directly
resolves this, but requires access to the source (available in
the Cargo registry cache at \texttt{\textasciitilde/.cargo/registry/src/}).

\paragraph{Verification.}
Five theorems in \texttt{Translation/NumIntegerSpec.lean}:

\begin{lstlisting}[caption={NI1 (symbolic): Integer::is\_even<u32> verified.},
                   label={lst:nispec}]
theorem num_is_even_symbolic (n : Nat) :
    evalFun [] NumIsEvenU32Fun 10 [.uint n] at () |=
      .pureOutput (· = .bool_ (n % 2 == 0)) :=
  ⟨.bool_ (n % 2 == 0), (), rfl, rfl⟩
\end{lstlisting}

The proof is by \texttt{rfl}: the kernel evaluates
\texttt{n \% 2 == 0} symbolically (the UB check `2 == 0' reduces to
\texttt{false} for the literal divisor, leaving only the arithmetic).
Ground instances (NI2--NI4) and a nocrash theorem (NI5) complete the set.

\subsection{TypeScript Cross-Language Unification}
\label{sec:eval-ts}

The TypeScript module (\texttt{TypeScript/ElabSpec.lean}, 498 lines) now
covers 14 functions and includes 6 unification theorems:

\begin{itemize}[leftmargin=2em]
  \item \textbf{T1--T14}: pure functions (return42, id, neg, add, sub, abs,
        max, isZero, mul, min, not\_gate, clamp) and the looping
        \texttt{sum\_to} function proved correct at four ground instances
        (n = 0, 1, 5, 10).
  \item \textbf{U1--U6}: cross-language unification theorems asserting that
        both a Rust LLBC implementation and a TypeScript implementation satisfy
        the \emph{same} \ProgramSpec predicate.
        Notable: U5 pairs \texttt{sumToFun} (Rust LLBC loop) with
        \texttt{tsSumToFun} (TypeScript \texttt{while} loop), both
        proving \texttt{sum\_to(10) = 45} from the identical specification.
\end{itemize}

The TypeScript \texttt{while} loop required extending the \texttt{TSStmt}
AST with a single mutable-update constructor (\texttt{set\_}), added in
3 lines to \texttt{AST.lean} and \texttt{Elaborator.lean}.
All existing proofs continue to build without modification.

\subsection{Lines of Code}

\begin{table}[t]
\centering
\caption{Lines of code by category (21 source files, 5\,361 total lines).}
\label{tab:loc}
\small
\begin{tabular}{@{}lrr@{}}
\toprule
\textbf{Category} & \textbf{Files} & \textbf{Lines} \\
\midrule
Foundation (Types, Monad, Ownership)                      & 3  &  281 \\
Specification (ProgramSpec, Satisfies, Examples)          & 3  &  427 \\
Translation (LLBC, Elaborator, ElabSpec, Semantics, Adequacy) & 5 & 2{,}095 \\
Loop invariants (LoopInvariant, FactInvariant, FibInvariant)  & 3 &   748 \\
Charon pipeline (CharonDefs, CharonSpec)                  & 2  &  756 \\
Tactic automation (VerifyFun, Examples)                   & 2  &  186 \\
TypeScript (AST, Elaborator, ElabSpec)                    & 3  &  868 \\
\midrule
\textbf{Total (Lean)}                                     & \textbf{21} & \textbf{5\,361} \\
Charon pipeline script (\texttt{charon2lean.py})          & 1  &  461 \\
\bottomrule
\end{tabular}
\end{table}

\subsection{Proof-to-Code Ratios}

\begin{table}[t]
\centering
\caption{Proof-to-code ratios. ``Def lines'' counts definition bodies;
         ``Proof lines'' counts proof bodies; ``Ratio'' = Proof/Def.}
\label{tab:ratios}
\small
\begin{tabular}{@{}lrrr@{}}
\toprule
\textbf{Module} & \textbf{Def lines} & \textbf{Proof lines} & \textbf{Ratio} \\
\midrule
\texttt{Spec/Satisfies.lean} (combinators)      &  51 & 136 & 2.7$\times$ \\
\texttt{Translation/ElabSpec.lean} (LLBC pure)  &  89 & 428 & 4.8$\times$ \\
\texttt{Translation/CharonSpec.lean} (Charon)   &  40 & 290 & 7.3$\times$ \\
\texttt{Tactic/Examples.lean} (macro proofs)    &  22 &  16 & 0.7$\times$ \\
\texttt{Translation/LoopInvariant.lean}         &  63 & 291 & 4.6$\times$ \\
\texttt{Translation/FibInvariant.lean}          &  47 & 186 & 4.0$\times$ \\
\texttt{TypeScript/ElabSpec.lean}               & 110 & 387 & 3.5$\times$ \\
\bottomrule
\end{tabular}
\end{table}

The \texttt{Tactic/Examples.lean} ratio of $0.7{\times}$ reflects the
proof automation macros: 16 theorems proved in 16 lines (one per theorem)
against 22 lines of statements. Without the macros, the same 16 theorems
would require approximately 176 lines ($11{\times}$ overhead).

The \texttt{CharonSpec.lean} ratio ($7.3{\times}$) is higher than the
hand-crafted \texttt{ElabSpec.lean} ($4.8{\times}$) because Charon-generated
functions use more temporaries (copy locals, overflow-check booleans), making
each function body larger relative to its proof.

\subsection{Verification Build Time}

We measured \lakeCmd{lake build Theorems} on an Apple M-series laptop
with 16\,GB RAM:

\begin{itemize}[leftmargin=2em]
  \item Full build from scratch (with \lakeCmd{lake exe cache get}): $\approx 4$ minutes.
  \item Incremental build (only lean-pl-verify changed): $\approx 30$ seconds.
  \item Mathlib compilation from source (no cache): 45--60 minutes.
\end{itemize}

The bulk of the build time is Mathlib compilation; lean-pl-verify itself
takes approximately 25--30 seconds to elaborate. \texttt{FibInvariant.lean}
takes the longest single file ($\approx 8$ seconds) due to the inductive
proof over $n.\texttt{toNat}$ iterations and the \texttt{native\_decide} call.

\subsection{Reproducibility}

The complete artifact is reproduced by:

\begin{lstlisting}[language={},caption={Reproduction commands.},label={lst:repro}]
$ git clone <repository> lean-pl-verify
$ cd lean-pl-verify
$ lake exe cache get      # download Mathlib cache (~1 GB, first time only)
$ lake build Theorems     # verifies all 258 theorems, 0 sorry
\end{lstlisting}

Successful completion of \lakeCmd{lake build Theorems} with no errors and no
sorry warnings is the reproducibility certificate.

\subsection{Comparison with Related Work}

\begin{table}[t]
\centering
\caption{Comparison with Aeneas on key dimensions.}
\label{tab:comparison2}
\small
\begin{tabular}{@{}lll@{}}
\toprule
\textbf{Dimension} & \textbf{Aeneas~\cite{ho2022aeneas}} & \textbf{This work} \\
\midrule
Embedding strategy    & Shallow (generated Lean functions) & Deep (inductive AST) \\
TCB                   & Charon + Aeneas translator         & Lean kernel only \\
Proof style           & Proofs over translated functions   & Proofs over AST \\
Loop invariants       & Yes (manual + loop summaries)      & Yes (manual, fuel-parametric) \\
Overflow model        & Explicit (Primitives library)      & Explicit (IntBounds checks) \\
TypeScript support    & No                                 & Yes (33 theorems) \\
Cross-language unif.  & No                                 & Yes (6 theorems) \\
Automated pipeline    & Yes (Charon integration)           & Yes (\S\ref{sec:eval-charon}) \\
Relational semantics  & No                                 & Yes (Semantics + Adequacy) \\
Sorry count           & N/A                                & \textbf{0 of 258 theorems} \\
Unsafe Rust           & No                                 & No \\
\bottomrule
\end{tabular}
\end{table}

The two systems are complementary rather than competing. Aeneas provides a
production-grade pipeline with a large body of verified cryptographic
code~\cite{aeneascrypto2026}. Lean-pl-verify provides a smaller,
kernel-minimal foundation that supports cross-language reasoning, a
relational semantics with adequacy, and an end-to-end Charon pipeline
to the same proof framework.
Combining the Aeneas translation pipeline with the \ProgramSpec unification
framework would be a natural direction for future work.
